\chapter{Boltzmann Weights, Partition Functions, and the Ideal Gas}

\lectureoverview{We derive the equilibrium probability distribution by
maximizing entropy at fixed normalization and mean energy.  The partition
function then yields the mean energy, entropy, and temperature by
differentiation.  The ideal gas supplies the first full example, followed by
a gravitational extension and a preview of energy fluctuations.}

\section{The Constrained Entropy Problem}

Begin with a large collection of \(\mathcal N\) identical, very weakly
interacting replicas of the same subsystem.  If \(n_i\) replicas occupy the
\(i\)-th state, then
\begin{equation}
 p_i=\frac{n_i}{\mathcal N},
 \qquad
 \sum_i p_i=1.
 \label{eq:sm4-occupation-probability}
\end{equation}
The multinomial counting of the preceding lecture gives the entropy per
replica,
\begin{equation}
 S[p]=-\sum_i p_i\log p_i.
 \label{eq:sm4-entropy}
\end{equation}
It is worth viewing this as an ordinary function of all the variables
\(p_1,p_2,\ldots\).  Its simple additive form is what makes the calculation
possible.

The replica collection has a fixed total energy \(\mathcal N E\), where
\(E\) is the mean energy per replica.  The second constraint is therefore
\begin{equation}
 \sum_i p_iE_i=E.
 \label{eq:sm4-energy-constraint}
\end{equation}
The levels \(E_i\) belong to the subsystem; the number \(E\) on the
right-hand side is their probability-weighted mean.  Maximizing
\eqref{eq:sm4-entropy} subject to
\eqref{eq:sm4-occupation-probability} and
\eqref{eq:sm4-energy-constraint} selects the overwhelmingly most probable
occupation fractions.

For convenience, minimize \(-S=\sum_i p_i\log p_i\).  Introduce one
Lagrange multiplier for each constraint:
\begin{equation}
 \begin{aligned}
 \Phi[p;\alpha,\beta]
 &=\sum_i p_i\log p_i\\
 &\quad+\alpha\left(\sum_i p_i-1\right)\\
 &\quad+\beta\left(\sum_i p_iE_i-E\right).
 \end{aligned}
 \label{eq:sm4-augmented-functional}
\end{equation}
The names \(\alpha\) and \(\beta\) are not yet physical assumptions.  At
this stage they are simply the two numbers required to enforce the two
constraints.

\begin{classroomqa}
\audiencequestion{Is the augmented function also a function of
\(\alpha\) and \(\beta\)?}
\lecturerresponse{Yes.  Varying it with respect to \(\alpha\) and \(\beta\)
restores the two constraints.  For the immediate calculation, however, we
hold those multipliers fixed and vary the probabilities \(p_i\).  Their
values are fixed afterward by normalization and the prescribed mean energy.
\lecturetimestamp{00:09:32--00:09:49}}
\end{classroomqa}

\begin{figure}[t]
\centering
\includegraphics[width=0.92\textwidth]{lecture_04_figure_02.png}
\caption{The augmented entropy functional after terms constant under
variation of the probabilities have been suppressed.}
\lectureframe{00:10:56}
\end{figure}

When differentiating with respect to a probability, the terms
\(-\alpha\) and \(-\beta E\) contribute nothing.  This is the abbreviated
functional visible on the board.  They remain part of
\eqref{eq:sm4-augmented-functional}; they have merely become irrelevant to
this particular derivative.

\section{The Boltzmann Distribution}

Choose any one probability.  The lecture chooses \(p_7\), partly because
\(1\) looks too special and \(7\) is a lucky number.  Nothing depends on
that choice.  The only calculus required is
\begin{equation}
 \frac{d}{dp_i}\bigl(p_i\log p_i\bigr)=\log p_i+1.
\end{equation}
Thus
\begin{equation}
 \frac{\partial\Phi}{\partial p_i}
 =\log p_i+1+\alpha+\beta E_i=0.
 \label{eq:sm4-stationarity}
\end{equation}
Solving for \(p_i\),
\begin{align}
 \log p_i&=-(1+\alpha)-\beta E_i,\\
 p_i&=e^{-(1+\alpha)}e^{-\beta E_i}.
\end{align}
The first factor is independent of the state label.  Write it as \(1/Z\):
\begin{equation}
 p_i=\frac{e^{-\beta E_i}}{Z}.
 \label{eq:sm4-boltzmann-weight}
\end{equation}

This is the first sudden payoff.  A routine constrained derivative has
produced an exponential law.  The relative probabilities obey
\begin{equation}
 \frac{p_i}{p_j}=e^{-\beta(E_i-E_j)}.
 \label{eq:sm4-relative-weight}
\end{equation}
The multiplier \(\beta\) controls how rapidly probability falls with
energy.  Large positive \(\beta\) concentrates the distribution near the
lowest levels; small positive \(\beta\) spreads appreciable weight farther
up the spectrum.  We can already see that changing \(\beta\) changes the
mean energy, although its identification with inverse temperature still
has to be proved.

\section{Normalization and the Partition Function}

Normalization determines the state-independent factor:
\begin{equation}
 1=\sum_i p_i
 =\frac{1}{Z}\sum_i e^{-\beta E_i},
\end{equation}
so
\begin{equation}
 Z(\beta)=\sum_i e^{-\beta E_i}.
 \label{eq:sm4-partition-function}
\end{equation}
This is the canonical partition function.  It first appeared as a
normalizing constant, but that description badly understates its role.
Once \(Z(\beta)\) is known, thermodynamic quantities follow from its
derivatives.

The energy levels in \eqref{eq:sm4-partition-function} are inputs.  They
come from the mechanics of the system together with fixed external
conditions: the volume and shape of the box, magnetic or electric fields,
and any other control parameters.  Changing a control parameter changes
the spectrum and therefore changes \(Z\).  For the present calculation
those parameters, and hence the set \(\{E_i\}\), are held fixed.

\begin{classroomqa}
\audiencequestion{What exactly do the axes of the exponential sketch
represent?  Is the vertical coordinate one particular \(p_i\)?}
\lecturerresponse{The plotted height is the probability assigned to a
level, while the horizontal ordering may be read either as the state label
\(i\) or, more physically, as the energy \(E_i\).  Every state has a fixed
energy.  Large \(\beta\) makes the probability fall steeply with that
energy; near \(\beta=0\), the weights are nearly flat over the accessible
levels.
\lecturetimestamp{00:19:39--00:20:09}}
\end{classroomqa}

\section{Mean Energy Hidden in the Partition Function}

The second constraint determines \(\beta\):
\begin{equation}
 E=\sum_i p_iE_i
 =\frac{1}{Z}\sum_i E_i e^{-\beta E_i}.
 \label{eq:sm4-mean-energy}
\end{equation}
Rather than evaluate this weighted sum afresh, differentiate
\eqref{eq:sm4-partition-function}:
\begin{equation}
 \frac{dZ}{d\beta}
 =-\sum_i E_i e^{-\beta E_i}.
\end{equation}
Comparison with \eqref{eq:sm4-mean-energy} gives
\begin{equation}
 E(\beta)
 =-\frac{1}{Z}\frac{dZ}{d\beta}
 =-\frac{d\log Z}{d\beta}.
 \label{eq:sm4-energy-logz}
\end{equation}
The logarithm is natural because its derivative supplies the factor
\(1/Z\).

\begin{figure}[t]
\centering
\includegraphics[width=0.82\textwidth]{lecture_04_figure_03.png}
\caption{Normalization, the mean-energy constraint, and the derivative of
the partition function appear together on the board.}
\lectureframe{00:26:24}
\end{figure}

\begin{classroomqa}
\audiencequestion{Is the lower board equation a new equation for \(Z\), or
is it obtained from the definition above it?}
\lecturerresponse{It is obtained by differentiating the definition
\(Z=\sum_i e^{-\beta E_i}\).  The derivative pulls down
\(-E_i\).  Dividing the result by \(Z\) then reproduces the weighted sum
that defines the mean energy.
\lecturetimestamp{00:25:58--00:26:37}}
\end{classroomqa}

\begin{classroomqa}
\audiencequestion{Does
\(dZ/d\beta=-E Z\) make \(Z\) an eigenfunction of differentiation?}
\lecturerresponse{Not in the usual sense.  An eigenvalue equation would
have a constant multiplying \(Z\), whereas here
\(E=E(\beta)\).  The equation tells us the mean energy as a function of
\(\beta\); it is not a fixed-eigenvalue problem.
\lecturetimestamp{00:28:49--00:29:43}}
\end{classroomqa}

The simplicity of \eqref{eq:sm4-energy-logz} is deceptive.  The spectrum
may be complicated, and calculating \(Z\) can be difficult, but once that
single function has been found the mean energy follows by one derivative.
This pattern recurs throughout statistical mechanics: elementary
operations applied in the right order yield results that would have been
hard to guess directly.

\begin{classroomqa}
\audiencequestion{Is \(E\) the total energy of all the replicas or the
average energy per box?}
\lecturerresponse{It is the average energy of one box.  The replica
collection has total energy \(\mathcal N E\).  Later, when the system itself
is chosen to be a whole box containing \(N\) gas particles, \(E\) will mean
the energy of that whole gas sample.
\lecturetimestamp{00:30:55--00:31:10}}
\end{classroomqa}

\begin{classroomqa}
\audiencequestion{Can we take the logarithm of
\(\sum_i e^{-\beta E_i}\) term by term?}
\lecturerresponse{No.  A logarithm turns a product into a sum, not a sum
into a sum of logarithms:
\[
 \log\!\left(\sum_i e^{-\beta E_i}\right)
 \ne \sum_i\log\!\left(e^{-\beta E_i}\right).
\]
This distinction becomes useful later, when independent subsystems make
their partition functions multiply.
\lecturetimestamp{00:31:10--00:31:58}}
\end{classroomqa}

\section{What Is Fixed, and What Is Allowed to Vary?}

There are several layers of variables in the same formulas, and the
classroom discussion deliberately separates them.

\begin{itemize}
 \item The probabilities \(p_i\) are the original unknowns.
 \item The levels \(E_i\) are fixed once the system and its control
 parameters are specified.
 \item The multiplier \(\alpha\) has been absorbed into the normalization
 \(Z=e^{1+\alpha}\).
 \item For a prescribed mean energy \(E\), the constraint fixes one number
 \(\beta\).  Across a family of equilibrium states, however, either
 \(E=E(\beta)\) or \(\beta=\beta(E)\) may be used.
\end{itemize}

\begin{classroomqa}
\audiencequestion{If the multipliers are numbers, why are \(\beta\) and
\(E\) now being treated as variables?}
\lecturerresponse{For one constrained problem they are indeed numbers.
But we may compare many equilibrium problems with the same spectrum and
different prescribed mean energies.  Along that family, \(\beta\) changes
with \(E\).  Fixing \(E\) determines \(\beta\); fixing \(\beta\) determines
\(E\).  These are two readings of the same relation.
\lecturetimestamp{00:31:58--00:33:27}}
\end{classroomqa}

\begin{classroomqa}
\audiencequestion{Has \(\alpha\) disappeared?}
\lecturerresponse{Only its name has disappeared.  The definition
\(e^{-(1+\alpha)}=1/Z\) absorbs it into \(Z\).  Keeping both \(\alpha\)
and \(Z\) would count the same normalization parameter twice.
\lecturetimestamp{00:33:29--00:33:49}}
\end{classroomqa}

Energy need not be the only conserved quantity.  If total momentum, particle
number, angular momentum, or another quantity is constrained, its
expectation value receives its own multiplier.  That enlargement changes
the ensemble but not the method.

\begin{classroomqa}
\audiencequestion{Does fixing only the mean energy presume too much about
the system?}
\lecturerresponse{It presumes that no other conserved quantity needs to be
conditioned on in the present problem.  If, for example, total momentum is
also fixed, add another constraint and another Lagrange multiplier.  By
contrast, volume, shape, and external fields are control parameters here:
they are held fixed and determine the energy levels themselves.
\lecturetimestamp{00:33:51--00:35:50}}
\end{classroomqa}

\begin{classroomqa}
\audiencequestion{Should \(\alpha\) appear as another argument of the
energy, and does \(Z\) depend on the levels \(E_i\)?}
\lecturerresponse{\(\alpha\) is already contained in \(Z\), so it is not
an independent argument.  The partition function certainly depends on the
fixed spectrum as well as on \(\beta\).  Changing a control parameter
changes the \(E_i\), and therefore changes \(Z\); during the present
\(\beta\)-derivatives the spectrum is held fixed.
\lecturetimestamp{00:35:52--00:37:10}}
\end{classroomqa}

\begin{classroomqa}
\audiencequestion{Are we varying the individual levels \(E_i\), or the
mean energy \(E\)?}
\lecturerresponse{The individual levels remain fixed.  The mean energy
changes because the probabilities change.  Thus the family is parameterized
by \(E\) or equivalently by \(\beta\), while the same list of levels
\(\{E_i\}\) is used throughout.
\lecturetimestamp{00:37:12--00:39:04}}
\end{classroomqa}

In practice temperature is often much easier to measure than the total
energy of a gas.  We can insert a thermometer; weighing the relativistic
mass difference and invoking \(E=mc^2\) would be correct but spectacularly
unhelpful.  To make \(\beta\) operational, we must now connect it to the
thermodynamic definition of temperature.

\section{Entropy Before Temperature}

Entropy logically comes first.  Insert the Boltzmann probabilities into
\eqref{eq:sm4-entropy}:
\begin{equation}
 \log p_i=-\beta E_i-\log Z.
\end{equation}
Then
\begin{align}
 S
 &=-\sum_i p_i\log p_i\\
 &=\beta\sum_i p_iE_i+\log Z\sum_i p_i\\
 &=\beta E+\log Z.
 \label{eq:sm4-entropy-logz}
\end{align}
The normalization sum supplies the final factor of one.  This is also the
first hint of the Helmholtz free energy, since later
\(F_{\mathrm H}=-T\log Z\) in units with \(k_B=1\).

\begin{classroomqa}
\audiencequestion{What happened to the factor \(1/Z\) in the first term?}
\lecturerresponse{It is already part of \(p_i\).  The sum
\(\sum_i p_iE_i\) includes the complete normalized probability and is, by
definition, the mean energy \(E\).
\lecturetimestamp{00:43:19--00:43:26}}
\end{classroomqa}

\begin{classroomqa}
\audiencequestion{Should there be a minus sign in
\(S=\beta E+\log Z\)?}
\lecturerresponse{No.  The entropy begins with a minus sign, while
\(\log p_i=-\beta E_i-\log Z\) contains two more minus signs.  They produce
the two positive terms displayed in \eqref{eq:sm4-entropy-logz}.
\lecturetimestamp{00:45:53--00:46:03}}
\end{classroomqa}

Temperature was defined thermodynamically by
\begin{equation}
 dE=T\,dS,
 \qquad
 \frac{1}{T}=\frac{dS}{dE},
 \label{eq:sm4-temperature-definition}
\end{equation}
with all other control parameters held fixed.  A first attempt to
differentiate \eqref{eq:sm4-entropy-logz} directly with respect to \(E\)
quickly becomes cluttered because \(\beta\) also depends on \(E\).  Taking
a total differential is cleaner:
\begin{align}
 dS
 &=\beta\,dE+E\,d\beta
   +\frac{d\log Z}{d\beta}\,d\beta\\
 &=\beta\,dE+E\,d\beta-E\,d\beta\\
 &=\beta\,dE.
 \label{eq:sm4-entropy-differential}
\end{align}
The cancellation uses \eqref{eq:sm4-energy-logz}.  Comparing
\eqref{eq:sm4-entropy-differential} with
\eqref{eq:sm4-temperature-definition} gives
\begin{equation}
 T=\frac{1}{\beta},
 \qquad
 \beta=\frac{1}{T}.
 \label{eq:sm4-beta-temperature}
\end{equation}

\begin{figure}[t]
\centering
\includegraphics[width=0.82\textwidth]{lecture_04_figure_04.png}
\caption{The two \(d\beta\) terms cancel, leaving
\(dS=\beta\,dE\).}
\lectureframe{00:50:48}
\end{figure}

\begin{classroomqa}
\audiencequestion{Does this cancellation incorrectly assume that
\(\log Z\) is independent of energy or temperature?}
\lecturerresponse{No.  For the fixed spectrum there is one independent
thermal parameter.  We may write \(Z=Z(\beta)\), then regard
\(\beta=\beta(E)\) along the equilibrium family.  The total differential
keeps the full \(d\beta\) dependence, and the identity
\(d\log Z/d\beta=-E\) is exactly what makes the cancellation work.
\lecturetimestamp{00:52:04--00:52:44}}
\end{classroomqa}

Throughout most of the calculation temperature is measured in energy units,
so \(k_B=1\).  Restoring laboratory units,
\begin{equation}
 \beta=\frac{1}{k_B T_{\mathrm K}},
 \label{eq:sm4-beta-laboratory}
\end{equation}
where \(T_{\mathrm K}\) is an absolute temperature in kelvin.  It is often
convenient to write \(T_{\mathrm e}=k_BT_{\mathrm K}\) and then use
\(\beta=1/T_{\mathrm e}\).

\begin{classroomqa}
\audiencequestion{For a volume of gas, is the system one molecule or the
whole collection?}
\lecturerresponse{That depends on how the system is defined.  Here the
system will be the entire box of gas, so \(E\) is the energy of all its
molecules.  For a sufficiently dilute ideal gas, one molecule can also be
treated as a subsystem because the particles do not interact; the
factorization below will show the relation between the two descriptions.
\lecturetimestamp{00:53:32--00:54:49}}
\end{classroomqa}

The canonical identities can now be collected:
\begin{align}
 p_i&=\frac{e^{-\beta E_i}}{Z},&
 Z(\beta)&=\sum_i e^{-\beta E_i},\label{eq:sm4-canonical-pair}\\
 E&=-\frac{d\log Z}{d\beta},&
 S&=\beta E+\log Z,&
 T&=\frac{1}{\beta}.
 \label{eq:sm4-canonical-identities}
\end{align}

\section{A Classical Ideal Gas}

Take \(N\) identical point particles of mass \(m\) in a box of volume
\(V\).  The gas is ideal when interactions between particles can be
neglected.  This may be a good approximation because the forces are weak
or because the gas is dilute enough that close encounters are rare.

A classical microstate is specified by all positions and momenta:
\begin{equation}
 (\bm q_1,\ldots,\bm q_N;\bm p_1,\ldots,\bm p_N).
\end{equation}
There are \(3N\) position coordinates and \(3N\) momentum coordinates, so
the state is a point in a \(6N\)-dimensional phase space.  With no
interactions, the Hamiltonian is
\begin{equation}
 H(\bm p)=\sum_{a=1}^{N}\frac{\bm p_a^2}{2m}
 =\sum_{\mu=1}^{3N}\frac{p_\mu^2}{2m}.
 \label{eq:sm4-ideal-hamiltonian}
\end{equation}

\begin{classroomqa}
\audiencequestion{Are there \(3N\) particles, and does summing over all
momentum components double-count the states?}
\lecturerresponse{There are \(N\) particles and \(3N\) momentum
components.  The sum in \eqref{eq:sm4-ideal-hamiltonian} contains one
quadratic term for each component.  The phase-space integral then integrates
each independent component exactly once, so there is no double counting
from that product.
\lecturetimestamp{01:02:29--01:04:04}}
\end{classroomqa}

For a continuous classical phase space, the state sum becomes
\begin{equation}
 Z_N
 =\frac{1}{N!\,h^{3N}}
 \int_{V^N}d^{3N}q
 \int_{\mathbb R^{3N}}d^{3N}p\,
 e^{-\beta H(\bm p)}.
 \label{eq:sm4-classical-partition}
\end{equation}
The lecture concentrates on the integrals and writes the same expression
without the phase-space cell \(h^{3N}\).  Restoring it makes \(Z_N\)
dimensionless and records the quantum origin of classical state counting.
Because \(h\) is independent of \(\beta\), it does not alter the energy
derivative used below.

The position integral is immediate.  The Hamiltonian is independent of
position, and each particle may lie anywhere in the box:
\begin{equation}
 \int_{V^N}d^{3N}q=V^N.
\end{equation}
The momentum integral factorizes into \(3N\) identical one-dimensional
Gaussian integrals:
\begin{equation}
 Z_N
 =\frac{V^N}{N!\,h^{3N}}
 \left[
  \int_{-\infty}^{\infty}dp\,
  e^{-\beta p^2/(2m)}
 \right]^{3N}.
 \label{eq:sm4-factorized-gas}
\end{equation}

The factor \(1/N!\) identifies configurations that differ only by a
permutation of identical particles.  For the mean energy at fixed \(N\)
and \(V\), it contributes only a \(\beta\)-independent constant to
\(\log Z_N\), so it disappears under
\(-\partial/\partial\beta\).  It is not disposable in every question:
together with \(h^{3N}\), it is essential for a dimensionless partition
function, an extensive entropy, and a consistent chemical potential.

Now set
\begin{equation}
 x=\sqrt{\frac{\beta}{2m}}\,p,
 \qquad
 dp=\sqrt{\frac{2m}{\beta}}\,dx.
\end{equation}
The Gaussian integral becomes
\begin{align}
 \int_{-\infty}^{\infty}dp\,e^{-\beta p^2/(2m)}
 &=
 \sqrt{\frac{2m}{\beta}}
 \int_{-\infty}^{\infty}dx\,e^{-x^2}\\
 &=\sqrt{\frac{2\pi m}{\beta}}.
 \label{eq:sm4-gaussian}
\end{align}

\begin{figure}[t]
\centering
\includegraphics[width=0.82\textwidth]{lecture_04_figure_05.png}
\caption{The momentum integral is reduced to the standard Gaussian by a
single rescaling.}
\lectureframe{01:13:12}
\end{figure}

Substitution into \eqref{eq:sm4-factorized-gas} gives
\begin{equation}
 Z_N
 =\frac{V^N}{N!\,h^{3N}}
 \left(\frac{2\pi m}{\beta}\right)^{3N/2}.
 \label{eq:sm4-ideal-partition}
\end{equation}

\begin{figure}[t]
\centering
\includegraphics[width=0.78\textwidth]{lecture_04_figure_06.png}
\caption{The completed ideal-gas partition function on the board.  The
phase-space cell \(h^{3N}\) is restored in
\eqref{eq:sm4-ideal-partition}.}
\lectureframe{01:14:38}
\end{figure}

It is useful to package the same result with the thermal de Broglie
wavelength
\begin{equation}
 \lambda_T=\frac{h}{\sqrt{2\pi mT}}
 \qquad (k_B=1).
\end{equation}
Then
\begin{equation}
 Z_N=\frac{1}{N!}\left(\frac{V}{\lambda_T^3}\right)^N.
 \label{eq:sm4-thermal-wavelength-form}
\end{equation}

For macroscopic \(N\), Stirling's formula gives
\begin{equation}
 N!\simeq N^Ne^{-N}\sqrt{2\pi N}.
\end{equation}
Keeping only the leading exponential dependence and defining
\(\rho=N/V\),
\begin{equation}
 \frac{V^N}{N!}
 \simeq
 \frac{1}{\sqrt{2\pi N}}
 \left(\frac{e}{\rho}\right)^N.
 \label{eq:sm4-density-rewrite}
\end{equation}
This is the lecture's neat density reorganization.  The square-root factor
is subleading in \(\log Z_N\), but it is displayed here rather than silently
dropped.

\begin{classroomqa}
\audiencequestion{If \(1/N!\) is independent of \(\beta\), does it make no
difference at all?}
\lecturerresponse{It makes no difference to the present mean-energy
derivative at fixed \(N\) and \(V\), because it adds a
\(\beta\)-independent constant to \(\log Z_N\).  It does matter for absolute
entropy and for variations in particle number.  The limited derivative
statement and the broader state-counting statement should not be confused.
\lecturetimestamp{01:27:01--01:27:51}}
\end{classroomqa}

\section{Energy and Equipartition}

Products are awkward to differentiate; logarithms turn them into sums.
From \eqref{eq:sm4-ideal-partition},
\begin{equation}
 \log Z_N
 =-\frac{3N}{2}\log\beta+\text{\(\beta\)-independent terms}.
 \label{eq:sm4-log-ideal-partition}
\end{equation}
Therefore
\begin{equation}
 E
 =-\frac{\partial\log Z_N}{\partial\beta}
 =\frac{3N}{2\beta}
 =\frac{3N}{2}T.
 \label{eq:sm4-ideal-energy}
\end{equation}
The energy per particle is
\begin{equation}
 \frac{E}{N}=\frac{3}{2}T
 \qquad (k_B=1),
\end{equation}
or, in laboratory units,
\begin{equation}
 \frac{E}{N}=\frac{3}{2}k_BT_{\mathrm K}.
 \label{eq:sm4-ideal-energy-lab}
\end{equation}

\begin{figure}[t]
\centering
\includegraphics[width=0.84\textwidth]{lecture_04_figure_07.png}
\caption{Only the \(-3N\log\beta/2\) term contributes to the ideal-gas
energy, giving \(E/N=3T/2\).}
\lectureframe{01:25:28}
\end{figure}

\begin{classroomqa}
\audiencequestion{Where does Boltzmann's constant enter the formula for
the energy?}
\lecturerresponse{The theorist's temperature \(T\) is an energy.  To use
kelvin, replace it by \(k_BT_{\mathrm K}\), giving
\(E/N=3k_BT_{\mathrm K}/2\).  The small value of \(k_B\) converts a
macroscopic-looking numerical temperature into the microscopic energy
scale of one molecule.
\lecturetimestamp{01:24:58--01:26:01}}
\end{classroomqa}

The three comes from the three directions of space.  The one-half comes
from the square-root dependence of each Gaussian integral on \(1/\beta\).
Each independent quadratic momentum coordinate contributes
\begin{equation}
 \frac{1}{2}T
\end{equation}
to the mean energy.  In \(d\) spatial dimensions the same argument gives
\(E/N=dT/2\).  This is the first appearance of equipartition.

The entropy calculation assigned at the end of the lecture follows from
\(S=\beta E+\log Z_N\).  Using Stirling's approximation,
\begin{equation}
 \frac{S}{N}
 \simeq
 \log\!\left[
  \frac{V}{N}
  \left(\frac{2\pi mT}{h^2}\right)^{3/2}
 \right]
 +\frac{5}{2},
 \label{eq:sm4-sackur-tetrode}
\end{equation}
up to the subleading finite-\(N\) correction.  Restoring \(k_B\) multiplies
the right-hand side by \(k_B\) and replaces \(T\) by
\(k_BT_{\mathrm K}\).  The result makes clear why the normalization factors
that dropped out of the energy are indispensable for the entropy.

\section{Putting the Gas in a Gravitational Field}

A question about placing the box on a neutron star introduces the
gravitational extension.  Before confronting a neutron star, consider a
uniform field \(g\) in an ordinary room.  Choose \(y=0\) at the bottom of a
box of base area \(A\) and height \(L\).  The Hamiltonian becomes
\begin{equation}
 H=\sum_{a=1}^{N}
 \left(\frac{\bm p_a^2}{2m}+mgy_a\right).
 \label{eq:sm4-gravity-hamiltonian}
\end{equation}
The horizontal position integrals still give \(A^N\), while each vertical
coordinate contributes
\begin{equation}
 I_y(L)
 =\int_0^Ldy\,e^{-\beta mgy}
 =\frac{1-e^{-\beta mgL}}{\beta mg}.
 \label{eq:sm4-finite-height-integral}
\end{equation}

\begin{classroomqa}
\audiencequestion{Should the vertical integral really run to infinity, or
only to the top of the box?}
\lecturerresponse{For a finite box it must run from \(0\) to its height
\(L\), as in \eqref{eq:sm4-finite-height-integral}.  Taking the upper limit
to infinity is a later simplification that models an atmosphere above a
surface rather than a finite room.
\lecturetimestamp{01:30:25--01:31:32}}
\end{classroomqa}

The full finite-height partition function is therefore
\begin{equation}
 Z_N(L)
 =
 \frac{1}{N!\,h^{3N}}
 \left[
  A\,
  \frac{1-e^{-\beta mgL}}{\beta mg}
  \left(\frac{2\pi m}{\beta}\right)^{3/2}
 \right]^N.
 \label{eq:sm4-gravity-partition-finite}
\end{equation}
The vertical factor depends on \(\beta\), so it must contribute to the
energy.  Indeed,
\begin{equation}
 \langle U\rangle
 =N\left[
  \frac{1}{\beta}
  -\frac{mgL}{e^{\beta mgL}-1}
 \right].
 \label{eq:sm4-finite-gravity-energy}
\end{equation}
This completes the small piece of integration that the lecture first leaves
as a nuisance calculation.

For an infinitely high atmosphere, \(L\to\infty\), and
\begin{equation}
 I_y(\infty)=\frac{1}{\beta mg}.
\end{equation}
The partition function becomes
\begin{equation}
 Z_N^{(\infty)}
 =
 \frac{1}{N!\,h^{3N}}
 \left[
  \frac{A}{\beta mg}
  \left(\frac{2\pi m}{\beta}\right)^{3/2}
 \right]^N.
 \label{eq:sm4-gravity-partition-infinite}
\end{equation}
Its \(\beta\)-dependent logarithm is
\begin{equation}
 \log Z_N^{(\infty)}
 =-\frac{3N}{2}\log\beta-N\log\beta+\mathrm{const}.
\end{equation}
The first term gives the kinetic energy \(3NT/2\); the second gives the
mean gravitational potential energy \(NT\).  Hence
\begin{equation}
 E_{\mathrm{tot}}=\frac{5N}{2}T
\end{equation}
for this idealized unbounded atmosphere.

\begin{figure}[t]
\centering
\includegraphics[width=0.83\textwidth]{lecture_04_figure_08.png}
\caption{The infinite-atmosphere partition function adds
\(-N\log\beta\) to \(\log Z\), the source of the mean potential energy.}
\lectureframe{01:36:38}
\end{figure}

\begin{classroomqa}
\audiencequestion{Would a very strong gravitational field require another
Lagrange multiplier?}
\lecturerresponse{Not merely because the field is strong.  Gravity changes
the Hamiltonian and therefore the energy levels or phase-space weight.  The
same energy multiplier \(\beta\) remains.  A neutron star introduces
additional physical complications, but the first modification is still to
put the gravitational potential energy into \(H\).
\lecturetimestamp{01:27:58--01:29:24}}
\end{classroomqa}

\section{Three Questions at the Close}

\begin{classroomqa}
\audiencequestion{Which equation is the Boltzmann distribution, and what
does it assert about equilibrium?}
\lecturerresponse{The distribution is
\[
 p_i=\frac{e^{-\beta E_i}}{Z}.
\]
It is the maximum-entropy distribution subject to normalization and a fixed
mean energy.  The mean energy must be computed from the same probabilities,
and normalization must still be enforced even though it looks almost
automatic in the final formula.
\lecturetimestamp{01:38:14--01:38:45}}
\end{classroomqa}

The procedure can feel mechanical: write a partition function, take a
logarithm, and differentiate.  Yet the resulting equations continue to be
surprising.  Familiarity teaches which simple operation should follow which,
but it does not make the subject less curious.

\begin{classroomqa}
\audiencequestion{Can these entropy ideas be applied to a collapsing star
or a black hole?  How can entropy rise while the matter occupies a smaller
region?}
\lecturerresponse{Stars are difficult mainly because their material details
are complicated.  A black hole has the cleaner benchmark: its
Bekenstein--Hawking entropy is the limiting gravitational entropy for the
given macroscopic constraints.  Collapse can increase entropy until that
black-hole value is reached; increasing it further requires changing the
available energy or size.
\lecturetimestamp{01:39:29--01:40:36}}
\end{classroomqa}

\begin{classroomqa}
\audiencequestion{Can the partition function also calculate fluctuations
in energy?}
\lecturerresponse{Yes.  One derivative gives the mean energy, and a second
derivative gives its variance:
\begin{align*}
 \frac{\partial^2\log Z}{\partial\beta^2}
 &=\langle E^2\rangle-\langle E\rangle^2\\
 &=\operatorname{Var}(E).
\end{align*}
In energy units this is related to the heat capacity by
\(\operatorname{Var}(E)=T^2C_V\).  In laboratory units,
\(\operatorname{Var}(E)=k_BT_{\mathrm K}^2C_V\).  This fluctuation relation
is the next calculation.
\lecturetimestamp{01:40:37--01:41:53}}
\end{classroomqa}

Fluctuations distinguish statistical mechanics from a purely continuous
thermodynamic description.  Their systematic study, associated especially
with Gibbs and Einstein, exposes the finite microscopic constituents behind
the smooth macroscopic laws.  That is where the next lecture begins.
